\documentclass[12pt]{article}
\usepackage[letterpaper,margin={1.5cm}]{geometry}
\usepackage{amsmath, amssymb, amsfonts}
\usepackage[utf8]{inputenc}
\usepackage[T1]{fontenc}
\usepackage[spanish]{babel}
\usepackage{tikz}
\usepackage{graphicx,enumitem}
\usepackage{multicol}
\setlength{\marginparsep}{12pt} \setlength{\marginparwidth}{0pt} \setlength{\headsep}{.8cm} \setlength{\headheight}{15pt} \setlength{\labelwidth}{0mm} \setlength{\parindent}{0mm} \renewcommand{\baselinestretch}{1.15} \setlength{\fboxsep}{5pt} \setlength{\parskip}{3mm} \setlength{\arraycolsep}{2pt}
\renewcommand{\sin}{\operatorname{sen}}
\newcommand{\N}{\ensuremath{\mathbb{N}}}
\newcommand{\Z}{\ensuremath{\mathbb{Z}}}
\newcommand{\Q}{\ensuremath{\mathbb{Q}}}
\newcommand{\R}{\ensuremath{\mathbb{R}}}
\newcommand{\C}{\ensuremath{\mathbb{C}}}
\newcommand{\I}{\ensuremath{\mathbb{I}}}
\graphicspath{{../imagenes/}{imagenes/}{..}}
\allowdisplaybreaks{}
\raggedbottom{}
\setlength{\topskip}{0pt plus 2pt}
\newcommand{\profesor}{Edward Y. Villegas}
\newcommand{\asignatura}{M\'ETODOS NUM\'ERICOS}
\newcommand{\diff}[3]{\frac{d^{#3} #1}{d#2^{#3}}}
\newcommand{\pdiff}[3]{\frac{\partial^{#3} #1}{\partial#2^{#3}}}
\newcommand{\abs}[1]{\left| #1 \right|}
\begin{document}
  \pagestyle{empty}
  \begin{minipage}{\linewidth}
    \centering
    \begin{tikzpicture}[very thick,font=\small]
      \node at (2,6) {\includegraphics[width=3.5cm]{logoudem}};
      \node at (9.5,6) {\includegraphics[width=9cm]{cbudem}};
      \node[fill=white,draw=white,inner sep=1mm] at (9.5,5.05) {\bf Permanencia con calidad, Acompa\~nar para exigir};
      \node[fill=white,draw=white,inner sep=1mm] at (7.5,4.2) {\Large\bf DEPARTAMENTO DE CIENCIAS B\'ASICAS};
      \draw (0,0) rectangle (18,3.5);
      \draw (0,2.5)--(18,2.5) (0,1.5)--(18,1.5) (15,4.2)--(18,4.2) node[below,pos=.5] {CALIFICACI\'ON} (15,2.5)--(15,7)--(18,7)--(18,3.5) (8.4,0)--(8.4,1.5) (15,0)--(15,1.5) (10,1.5)--(10,2.5);
      \node[right] at (0,3.2) {\bf Alumno:}; \node[right] at (15,3.2) {\bf Carn\'e:};
      \node[right] at (0,2.2) {\bf Asignatura:};
      \node at (6,1.95) {\asignatura};
      \node[right] at (10,2.2) {\bf Profesor:};
      \node at (15,1.95) {\profesor};
      \node[right] at (0,1.2) {\bf Examen:};
      \draw (3.8,.9) rectangle (4.4,1.3); \node[left] at (3.8,1.1) {Parcial:};
      \draw (3.8,.2) rectangle (4.4,.6); \node[left] at (3.8,.4) {Previa:};
      \draw (7.4,.9) rectangle (8,1.3); \node[left] at (7.4,1.1) {Final:};
      \draw (7.4,.2) rectangle (8,.6); \node[left] at (7.4,.4) {Habilitaci\'on:};
       \node at (4, .4) {X}; % Quiz
      \node[right] at (10,.5) {3}; % Número de grupo
      \node[right] at (10,1.) {14 de abril de 2016}; % Fecha de presentación
      \node[right] at (8.4,1.05) {\bf Fecha:}; \node[right] at (8.4,.45) {\bf Grupo:};
      \node[align=center,text width=3cm,font=\footnotesize] at (16.5,.75) {\centering\bf Use solo tinta\\y escriba claro};
    \end{tikzpicture}
  \end{minipage}
Para el desarrollo de los cálculos puede usar \textbf{exclusivamente calculadora} (no se permite el uso de portátil o celulares en el examen), y % {presencial}
puede disponer de sus apuntes de clase. % {Quiz}
Todo valor reportado debe aproximarse a \textbf{5 CIFRAS SIGNIFICATIVAS} con \textbf{REDONDEO SIMÉTRICO} (no es necesario en enteros y valores dados en enunciado)% {Todos}
. Recuerde el uso del separador decimal acorde a la \textbf{NTC} % {Todos}
.
Indique \textbf{clara y explícitamente la respuesta final} de cada pregunta % {Todos}
en \textbf{lapicero} (requisito en caso de reclamación) % {presencial}
, y \textbf{justifique todas sus respuestas y procedimientos. Si algún elemento solicitado, en teoría no puede realizarse, indíquelo y justifique por que no se puede realizar lo solicitado. % {Todos}
}.
\vspace{-.5cm}
  \begin{enumerate}[leftmargin=*,widest=9]
     \item \(2.0\) Dada la función \(\exp(-(x-2)^2)\sin(x-2+\pi/2)\), aproxime \(f^\prime (2)\) con \(h=0.2\).

\textbf{R/} Al ser una función analítica y continua alrededor del punto de interés (de hecho en todos los reales), podemos usar puntos adelante y atrás del solicitado, y así aplicar un esquema central de 3 puntos para la primera derivada.
    \[
    f^\prime(2) = \frac{f(2.2000) - f(1.8000)}{2(0.2)} = \frac{0.94164 - 0.94164}{0.40000} = 0
    \]
    \item \(2.0\) Dados los puntos que representan una función continua
    \[
    \begin{array}{lr}
    x & f(x)\\
    4 & 8\\
    6 & 9\\
    8 & 9\\
    10 & 10\\
    12 & 8    
    \end{array}
    \]
    Determine el valor de la integral \( \int^{12}_4 f(x)dx\), usando criterio de mayor precisión con el mismo método.

\textbf{R/} Primero, nótese que la separación entre los puntos es la misma, por lo cual aplican los métodos para puntos equidistantes (\(h=2\)). La mayor precisión siempre corresponde al uso de la aproximación de mayor grado. En este caso, al tratarse del uso de una forma con repetición del mismo método se debe seleccionar aquel de mayor grado que cumpla con la restricción del número de intervalos. En este caso, observamos 4 intervalos, por lo que el método de mayor precisión con repetición corresponde a Simpson \(1/3\), por ser múltiplo de 2 pero no de 3 (para \(3/8\) se requiere intervalos múltiplos de 3).
\[
\int^{12}_4 f(x)dx = \frac{2}{3}\left(f(4) + f(12) + 2 f(8) + 4(f(6) + f(10) \right) = \frac{2}{3}\left(8 + 8 + 2 (9) + 4((9) + (10) \right) = 73.333
\]
   \item (\(1.0\)) ¿Es posible usar solo octágonos para mallar un dominio bidimensional para integración sobre elementos finitos? Justifique.

\textbf{R/} No es posible usar solo octágonos. Para los mallados en problemas de integración sobre elementos finitos se usan triángulos (generalmente), cuadriláteros y hexágonos (este último poco común) por su propiedad de ser formas teselables, que corresponde que simétricamente llenan con facilidad un dominio arbitrario y a que llevan a formas con propiedades idóneas en las funciones interpolantes.
  \end{enumerate}
\end{document}
